\chapter{Optional: sandboxes for builds and tests}\label{ch:sandboxes}

\emph{You want the builder of each run to build and test your project in a container with the
right toolchain and, say, its own database, instead of on the server itself.} That is what a
\term{sandbox} is. Sandboxes are off after installing Agern; skip this chapter if your projects
build fine with what the server has.

\section{What a sandbox is, and is not}\label{sec:sandbox-what}

A project describes its container in a file in its repository, \file{.agern/sandbox.yml}, the
\term{sandbox definition}: a toolchain image, an optional setup script, sidecar services such as
PostgreSQL, cache directories and limits, and optionally further images to test on, called
\term[platform]{platforms}. Once an admin has approved exactly that version of the file, each run
of the project gets its own container. The builder edits the files on the server and runs builds
and tests in the container, which sees the run's working copy at \file{/workspace} and nothing
else from the server. The containers run in rootless Podman as the \texttt{agern} user.

A sandbox is a toolchain and resource boundary, \emph{not} a security boundary: the agent itself
still runs on the server as \texttt{agern} (chapter~\ref{ch:what}). Never put secrets into the
file; everybody who can see the project can read it.

\section{Turn sandboxes on (on the server)}\label{sec:sandbox-enable}

\begin{enumerate}
  \item Install rootless Podman and its helpers. The package only suggests them, and
    \cmd{agern sandbox enable} never installs packages itself:
\begin{lstlisting}
sudo apt install podman uidmap netavark aardvark-dns crun slirp4netns
\end{lstlisting}
  \item Turn sandboxes on:
\begin{lstlisting}
sudo agern sandbox enable
\end{lstlisting}
    It gives the \texttt{agern} user a range of subordinate user and group ids, starts the sandbox
    service, and runs a self-test exactly as the server will use Podman (it pulls the small image
    \texttt{docker.io/library/busybox}, checks networking, limits and file ownership). Only when
    everything passes does it set \texttt{FACTORY\_SANDBOX=podman} in \file{/etc/agern/agern.env}
    and restart Agern, and only if nobody's work is in progress; otherwise it tells you to run
    \cmd{sudo agern restart} later.
  \item Check the result:
\begin{lstlisting}
sudo agern sandbox status
sudo agern doctor
\end{lstlisting}
\end{enumerate}

Podman needs cgroups v2 and unprivileged user namespaces, the defaults on Debian~13 and
Ubuntu~24.04. Limits, disk use and troubleshooting are in
\repodoc{INSTALL-LINUX.md\#13-sandboxes}{\texttt{INSTALL-LINUX.md}, section~13}.

\section{Describe the sandbox (in the project)}\label{sec:sandbox-file}

\begin{enumerate}
  \item Add \file{.agern/sandbox.yml} to the project's default branch. Only \texttt{image} is
    required, and it must be fully qualified and come from an allowed registry. A minimal file for a
    Gradle project with a database:
\begin{lstlisting}[style=conf]
image: docker.io/library/eclipse-temurin:21-jdk
services:
  postgres:
    image: docker.io/library/postgres:17
    env: { POSTGRES_DB: app, POSTGRES_USER: app, POSTGRES_PASSWORD: app }
    ready: { command: "pg_isready -h 127.0.0.1 -U app -d app" }
caches: [ /root/.gradle ]
\end{lstlisting}
    The database is reachable from the sandbox under the host name \texttt{postgres}. Every key,
    platforms, a non-root user and examples for other toolchains are in
    \repodoc{SANDBOXES.md}{\texttt{docs/SANDBOXES.md}}.
  \item Merge it into the default branch. Runs read the file from the commit they start from.
\end{enumerate}

\section{Approve it (as an admin)}\label{sec:sandbox-approve}

A new or changed definition waits for an admin. Until then, runs build and test on the server as
before.

\begin{enumerate}
  \item On the \ui{Admin} page, the \ui{Sandboxes} card lists the running sandboxes and, under
    \ui{Waiting for approval}, the projects whose file is new (\ui{needs approval}) or changed
    since its approval (\ui{changed}) (figure~\ref{fig:admin-sandboxes}). The project page's
    \ui{Sandbox} section links to the same page.
  \item Open the project's card. The \ui{Sandbox} page shows what the file asks for: images, the
    setup script, services, caches and resources (figure~\ref{fig:sandbox-pending}).
  \item Approve only what you would run on the server yourself. Choose \ui{Approve\ldots}, read the
    summary and choose \ui{Approve this version}. The page now shows the approved definition and
    who approved it (figure~\ref{fig:sandbox-approved}); \ui{Revoke\ldots} withdraws the approval.
\end{enumerate}

Every change of the file needs a new approval. A run that changes the file itself still uses the
version it started from, so a change can never approve itself.

\screenshot{44-admin-sandboxes-desktop}{The Sandboxes card on the Admin page: running sandboxes and definitions waiting for approval.}{admin-sandboxes}

\screenshot[0.85\linewidth]{43-sandbox-pending-desktop}{A sandbox definition waiting for an admin's approval.}{sandbox-pending}

\screenshot[0.85\linewidth]{42-sandbox-approved-desktop}{An approved sandbox: its definition, who approved it, and the file.}{sandbox-approved}

\section{What a run shows}\label{sec:sandbox-run}

The run page names the sandbox in its \ui{SANDBOX} line: the image, the version of the
definition and the platforms (you can see it in figure~\ref{fig:run-started} in
chapter~\ref{ch:run}). Without an approved definition the line says why, for example that its
sandbox definition is not approved and the run builds and tests on the host. The builder's
transcript shows each command it runs in the sandbox, with its output.

The spec writer plans the sandbox's use in a \ui{Verification} section of the spec. The ux and
ui-design reviewers start the app in a separate \term{review sandbox}: a container from the same
definition, with its own services and the working copy mounted read-only, shared by both reviewers
and removed when their review round ends. When they used it, the \ui{SANDBOX} line says
\ui{review sandbox \checkmark}, and their transcripts show each command marked \ui{review}. When all
sandbox slots of the server are in use, the reviewers simply review on the server as before.

\whatyouhave
Sandboxes are on, your project describes its container, and an admin has approved it: the next
run builds and tests in its own container. Next, take a tour of a project in
chapter~\ref{ch:project}.
